\appendix

\section{Key Algorithms}
\label{appendix:simulator_core}

\vheading{Simulator}
The core logic of the simulator is summarized in Algorithm \ref{alg:simulator}. At a high-level, the simulator handles discrete events and increments a clock on the completion of each event. Events are classified into different types based on their effect --- for instance, some events indicate the enqueuing of a kernel on the GPU, while others indicate a synchronization between device streams.

The scheduler (Algorithm \ref{alg:scheduler}) handles updating the state of each host and device in the simulated cluster and adding or removing operations from the corresponding queues. Resources that are busy will cause any new ops targeting them to be queued, effectively simulating blocking delays by deferring their execution. Every time an operation completes execution, the next scheduler tick pops an operation from the resource-specific queue and adds an $\texttt{EndEvent}$ marking its completion in the future. Every discrete $\texttt{EndEvent}$ in the top-level queue is followed by a scheduler tick, ensuring that operations do not block forever on resources.
\begin{algorithm}
\caption{Core Discrete-Event Simulator Algorithm}
\label{alg:simulator}
\begin{algorithmic}[1]
\Procedure{Simulate}{config, host\_op\_trace}
    \State $\textit{time} \gets 0$
    \State $\textit{event\_queue} \gets \text{ PriorityQueue()}$
    \State $\textit{cluster} \gets \text{ Cluster(config)}$
    \State $\textit{scheduler} \gets \text{ Scheduler(config, cluster)}$
    
    \Statex \LComment{Init queue with host ops and inter-host-op overheads from the input trace}
    \For{each $\textit{host\_op} \lor \textit{overhead}$  in $\textit{host\_op\_trace}$}
        \State $\textit{event} \gets \text{ HostOpArrivalEvent}(\textit{host\_op})$
        \State $\textit{overhead} \gets \text{ HostOverhead}(\textit{overhead\_dur})$
        \State $\textit{event\_queue.put(event)}$
        \State $\textit{event\_queue.put(overhead)}$

    \EndFor
    
    \Statex
    \While{$\neg \textit{event\_queue.empty()}$}
        \LComment{Get the next chronological event}
        \State $\textit{event} \gets \textit{event\_queue.get()}$
        
        \Statex \LComment{Update simulation time.}
        \State $\textit{time} \gets \textit{event.end\_time}$
        
        \Statex \Comment{Handle event polymorphically based on its type.}
        \State $\textit{new\_events} \gets \textit{event.handle\_event(scheduler)}$
        
        \LComment{Add newly generated events to the queue}
        \For{each $\textit{new\_event}$ in $\textit{new\_events}$}
            \State $\textit{event\_queue.put(new\_event)}$
        \EndFor
        
    \EndWhile
    
    \Statex
    \Return $\textit{time}$
\EndProcedure
\end{algorithmic}
\end{algorithm}

\begin{algorithm}
\caption{Scheduler Event Handling Logic}
\label{alg:scheduler}
\begin{algorithmic}[1]
\Procedure{Event.handle\_event}{scheduler}
    \State \Comment{The logic here is polymorphic, depending on the concrete event type.}
    \Statex
    
    \If{$\textit{event}$ is an \textbf{OpArrivalEvent}}
        \LComment{An operation from the trace has arrived (e.g., kernel launch).}
        \State \Return $\textit{scheduler.schedule\_operation(event.op)}$
    
    \ElsIf{$\textit{event}$ is an \textbf{EndEvent}}
        \LComment{An operation has finished, freeing a resource.}
        \State \Return $\textit{scheduler.op\_complete(event.op)}$
    \ElsIf{$\textit{event}$ is a \textbf{ScheduleEvent}}
        \LComment{A global scheduling tick occurs.}
        \State $\textit{newly\_started\_ops} \gets \textit{scheduler.schedule()}$
        \State \Return $\text{create\_end\_events}(\textit{newly\_started\_ops})$
        
    \Else
        \LComment{Handle other event types (e.g., sync, collective).}
        \State \Return $\textit{handle\_other\_events(scheduler, event)}$
    \EndIf
\EndProcedure
\end{algorithmic}

\vspace{0.5cm}

\begin{algorithmic}[1]
\Procedure{Scheduler.schedule()}{}
\State $\textit{newly\_started\_ops}$ = $\emptyset$
    \For{each $\textit{device, stream}$ in $\textit{scheduler.cluster}$ }
        \If{$\textit{device.is\_busy()} \lor \textit{stream.is\_busy()}$}
            \LComment{A required resource is busy, so don't deque}
            \State $\textbf{continue}$
        \Else
            \LComment{Resources are free, so process the op}
            \State $\textit{op} \gets \textit{device\_queue}$.front() 
            \State $\textit{device.set\_busy()}$
            \State $\textit{stream.set\_busy()}$
            
            \State $\textit{duration} \gets \text{get\_runtime(op)}$
            \State $\textit{end\_time} \gets \textit{current\_time} + \textit{duration}$
            \State $\textit{end\_event} \gets \text{ EndEvent}(\textit{end\_time}, \textit{op})$
        \EndIf
    \EndFor
    \State \Return $\textit{newly\_started\_ops}$
\EndProcedure
\end{algorithmic}

\vspace{0.5cm}

\begin{algorithmic}[1]
\Procedure{Scheduler.op\_complete}{completed\_op}
    \State $\textit{device}, \textit{stream} \gets \text{completed\_op.get\_resources()}$
    
    \LComment{Mark the resources as free}
    \State $\textit{device.set\_free()}$
    \State $\textit{stream.set\_free()}$
    
    \If{$\textit{device.wait\_queue.is\_not\_empty()}$}
        \State \Comment{Check for and schedule the next pending operation}
        \State $\textit{next\_op} \gets \textit{device.wait\_queue.get\_next()}$
        \State \Return $\text{schedule\_operation}(\textit{next\_op})$
    \Else
        \State \Return $\emptyset$ \Comment{No pending work for this resource}
    \EndIf
\EndProcedure
\end{algorithmic}
\end{algorithm}



\begin{algorithm}
\caption{Synchronization Wait Map Structures}
\label{alg:waitmaps}
\begin{algorithmic}[1]
\algblock{CudaEventWaitMap}{EndCudaEventWaitMap}
\CudaEventWaitMap
    \State \Comment{\textbf{Structure:} Map from a CUDA (event ID, version) pair to a list ops waiting for it. Versions track re-use of the same CUDA event handle.}
    \State $\textit{events}: (\textit{event\_id}, \textit{version}) \to \textit{waiting\_ops}$
    \Statex
    \Procedure{BlockOnEvent}{event\_id, version, op}
        \LComment{An operation `op` (from a host or stream) blocks on a future event.}
        \State $\textit{events}[\textit{event\_id},\textit{version}].add(op)$
        \State Stall the host/stream associated with $\textit{op}$
    \EndProcedure
    \Statex
    \Procedure{ReleaseWaiters}{event\_id, version}
        \State \Comment{The event has been recorded; release all waiting operations for scheduling.}
        \State $\textit{released\_ops} \gets \textit{events}.pop(\textit{event\_id},\textit{version})$
        \For{each $op$ in $released\_ops$}
        \State Free the host/stream associated with $op$ by creating the associated $EndEvent$ instances
        \EndFor
        \State \Return $\textit{released\_ops}$
    \EndProcedure
\EndCudaEventWaitMap
\end{algorithmic}

\vspace{0.5cm}
\hrule
\vspace{0.5cm}
\begin{algorithmic}[1]
\algblock{NetworkCollectiveWaitMap}{EndNetworkCollectiveWaitMap}
\NetworkCollectiveWaitMap
    \State \Comment{\textbf{Structure:} A map from a NCCL collective's unique ID to its list of participant kernels.}
    \State $\textit{collectives}: (\textit{nccl\_group\_id}, \textit{call\_idx}) \to \textit{kernels}$
    \Statex
    \Procedure{JoinCollective}{kernel}
        \State \Comment{A device's kernel joins a collective operation and waits for peers.}
        \State $collectives[\textit{group\_id}, \textit{call\_idx}].add(kernel)$
        \State $\textit{wait\_list} \gets collectives[\textit{group\_id}, \textit{call\_idx}]$
        \Statex
        \If{$\text{length}(\textit{wait\_list}) = \textit{kernel.num\_ranks}$}
            \LComment{The last worker has arrived; the collective can proceed.}
            \State $collectives.pop(\textit{group\_id}, \textit{call\_idx})$
            \LComment{Return all kernels to be scheduled.}
            \State \Return $\textit{wait\_list}$ 
        \Else
            \State \Comment{Not all workers have arrived; keep blocking.}
            \State \Return $\emptyset$
        \EndIf
    \EndProcedure
\EndNetworkCollectiveWaitMap
\end{algorithmic}
\end{algorithm}

\vfill\break
The simulator maintains two global structures to track synchronization across hosts and accelerators --- the CUDA Event Wait Map, and the Network Collective Wait Map. These are detailed in Algorithm \ref{alg:waitmaps}. Corresponding event types that perform synchronization lookup these maps in their handlers. For example, in the case of $\texttt{cudaEventSynchronize}$, the specific device stream or host that intends to block while waiting for a CUDA event inserts an entry in the wait-map (keeping its associated resources from processing ops), while the matching $\texttt{cudaEventRecord}$ triggers the resources to be freed by creating $\texttt{EndEvent}$ instances for the blocking ops.

A similar mechanism is used to model collectives --- each worker makes an entry in the collective wait map. Once the final worker of the collective joins, all the corresponding streams are unblocked and can proceed. In this case, $\texttt{EndEvent}$s are scheduled with a timestamp after the predicted duration of the collective; effectively, \sysname models the delays involved in starting a collective using a global sync point and then assumes that workers move in lockstep. Any effects associated with the on-the-wire time of the collective can thus be abstracted away in the predicted time --- while this is not completely faithful to the setup/teardown of NCCL collectives, it is sufficiently accurate for an end-to-end accounting of latency.

These fairly simple structures can express a wide variety of possible synchronization behaviors since they operate at the CUDA stream level. Computation streams can overlap with collectives since each stream is a separate resource. The host queue can block on a specific CUDA event or a device stream, deferring the execution of future CUDA API calls. An arbitrary pipeline parallel schedule is a combination of such synchronizing events, and thus \sysname can trivially capture these behaviors without any explicit modeling.


\section{Per-kernel prediction accuracy}
\label{appendix:predictions}
The default predictors in \sysname use random forest regressors trained on per-kernel runtime data. Tables \ref{tab:h100_kernel_accuracy}, \ref{tab:v100_kernel_accuracy} and \ref{tab:a40_kernel_accuracy} include metrics on the prediction error of the kernels trained for Megatron-LM (H100, V100) and PyTorch FSDP (A40). All results use a random 80:20 training/test data split.

As a general theme, we observe the same characteristics as \cite{geoffrey2021habitat}, \cite{zhu2020daydream} --- a small portion of the kernels are responsible for a significant portion of end-to-end prediction error (matmuls for language models, convolution kernels for vision models). As a result, even large percentage-wise errors in several other kernels do not cause any significant degradation in end-to-end accuracy.

In keeping with this observation, we conduct more extensive profiling of these heavy-hitter kernels --- sweeping a large space of input dimensions for convolution/matmul. The remaining kernels are scraped from traces, collecting by running a single-layer LLaMa/OPT/vision model over a range of batch sizes and tensor-parallel dimensions (since other optimizations like pipeline parallelism do not affect the runtime of a single kernel). The training set for the heavy-hitter kernels included $\approx$ 42k individual points, compared to a few thousand points each for the rest.

\begin{table}[htbp]
\centering
\begin{tabular}{l l}
\toprule
\textbf{Kernel} & \textbf{MAPE}\\
\midrule
RadixSortOnesweepKernel & 7.80\% \\
cuComputeGradGammaBeta & 7.95\% \\
masked\_softmax\_warp\_backward & 0.73\% \\
compute\_num\_of\_partial\_segments & 7.37\% \\
unrolled\_elementwise\_kernel & 5.80\% \\
write\_num\_of\_segments & 7.27\% \\
cuApplyLayerNorm & 1.98\% \\
MemcpyHtoD & 14.23\% \\
CatArrayBatchedCopy\_aligned16\_contig & 5.79\% \\
cuComputeGradInput & 3.50\% \\
MemcpyDtoH & 7.85\% \\
compute\_grad\_weight & 3.63\% \\
at\_cuda\_detailcubDeviceScanKernel & 5.37\% \\
cublasSgemm\_v2 & 3.65\% \\
cublasSgemmStridedBatched & 2.22\% \\
indexSelectLargeIndex & 1.88\% \\
multi\_tensor\_apply\_kernel & 1.68\% \\
at\_cuda\_detailcubDeviceScanInitKernel & 6.99\% \\
triu\_tril\_kernel & 4.38\% \\
vectorized\_elementwise\_kernel & 8.44\% \\
krn\_partial\_segment\_offset & 55.35\% \\
RadixSortExclusiveSumKernel & 14.30\% \\
CatArrayBatchedCopy & 43.71\% \\
fused\_dropout\_kernel\_vec & 1.50\% \\
index\_elementwise\_kernel & 12.86\% \\
sum\_and\_scatter & 48.82\% \\
MemcpyDtoD & 0.00\% \\
reduce\_kernel & 16.75\% \\
RadixSortHistogramKernel & 9.01\% \\
masked\_softmax\_warp\_forward & 1.00\% \\
cuComputePartGradGammaBeta & 4.12\% \\
krn\_partials\_per\_segment & 7.16\% \\
elementwise\_kernel & 10.28\% \\
elementwise\_kernel\_with\_index & 24.67\% \\
thrustcuda\_cubcore\_kernel\_agent & 12.51\% \\
Memset & 13.25\% \\
\bottomrule
\end{tabular}
\caption{Mean absolute percentage error on a held-out validation set, trained on H100 kernel runtimes. Important kernel types for Megatron-LM models include \texttt{cublasSgemm\_v2} and \texttt{cublasSgemmStridedBatched}, where we have <5\% prediction error. Kernels with large percentage-wise errors are extremely short in duration, and thus do not impact end-to-end latency significantly.}
\label{tab:h100_kernel_accuracy}
\end{table}

\begin{table}[htbp]
\centering
\begin{tabular}{l l}
\toprule
\textbf{Kernel} & \textbf{MAPE}\\
\midrule
scaled\_masked\_softmax\_warp\_backward & 0.41\% \\
at\_cuda\_detailcubDeviceScanKernel & 5.80\% \\
sum\_and\_scatter & 49.87\% \\
write\_num\_of\_segments & 30.59\% \\
vectorized\_elementwise\_kernel & 11.44\% \\
indexSelectLargeIndex & 7.20\% \\
elementwise\_kernel & 26.48\% \\
krn\_partial\_segment\_offset & 48.18\% \\
fused\_dropout\_kernel\_vec & 1.03\% \\
index\_elementwise\_kernel & 10.47\% \\
cuApplyLayerNorm & 1.36\% \\
elementwise\_kernel\_with\_index & 31.91\% \\
compute\_num\_of\_partial\_segments & 10.39\% \\
cuComputePartGradGammaBeta & 3.05\% \\
MemcpyDtoH & 39.56\% \\
unrolled\_elementwise\_kernel & 13.89\% \\
Memset & 36.75\% \\
MemcpyHtoD & 25.61\% \\
CatArrayBatchedCopy & 105.45\% \\
krn\_partials\_per\_segment & 11.16\% \\
cuComputeGradInput & 1.80\% \\
cublasSgemm\_v2 & 4.58\% \\
compute\_grad\_weight & 2.23\% \\
triu\_tril\_kernel & 11.76\% \\
multi\_tensor\_apply\_kernel & 3.40\% \\
RadixSortHistogramKernel & 9.00\% \\
at\_cuda\_detailcubDeviceScanInitKernel & 14.83\% \\
masked\_softmax\_warp\_forward & 1.20\% \\
RadixSortExclusiveSumKernel & 38.65\% \\
thrustcuda\_cubcore\_kernel\_agent & 32.53\% \\
cublasSgemmStridedBatched & 1.84\% \\
cuComputeGradGammaBeta & 18.95\% \\
CatArrayBatchedCopy\_aligned16\_contig & 20.19\% \\
reduce\_kernel & 24.64\% \\
RadixSortOnesweepKernel & 13.54\% \\
MemcpyDtoD & 33.25\% \\
scaled\_masked\_softmax\_warp\_forward & 0.48\% \\
softmax\_warp\_backward & 1.04\% \\
\bottomrule
\end{tabular}
\caption{Mean absolute percentage error on a held-out validation set, trained on V100 kernel runtimes. Important kernel types for Megatron-LM models include \texttt{cublasSgemm\_v2} and \texttt{cublasSgemmStridedBatched}, where we have <5\% prediction error. Kernels with large percentage-wise errors are extremely short in duration, and thus do not impact end-to-end latency significantly.}
\label{tab:v100_kernel_accuracy}
\end{table}


In contrast to computation operations, there is a much smaller set of network collectives (<10) that is used in deep learning workloads. Furthermore, the input space of these operators typically much smaller, typically comprising only two parameters -- number of workers and input size. This allows us to devise a simple policy for modeling these operations. We first  collect performance data in a fashion similar to \texttt{nccl-tests}. We only sample data in the range that is generally relevant for training workloads ranging from tens of megabytes to tens of gigabytes. We then use our regression pipeline to interpolate within this range. While this affects generalization to dimensions outside the range of the training set, this does not pose a problem in practice since the collective sizes are bounded by the batch size, model parameters and accelerator memory. %


Automatically generated fused kernels pose a unique challenge due to an explosion in generated kernel signatures --- arising from a large number of op combinations. We address this by collecting information from the compiler IR about the content of the kernels rather than just their inputs. In our experiments with the compiler-fused Triton kernels used in \texttt{torch.compile}, features such as the number of primitive Triton language instructions (add, sub etc.) in the kernel definition proved valuable in predicting kernel runtimes. The corresponding training data was collected by sweeping workload traces for different models/batch sizes and extracting the relevant features/runtimes. Through this approach, we achieve comparable accuracy to that of the kernel predictions trained on focused micro-benchmarks. 

\vspace{1em}

\begin{table}[htbp]
\centering
\begin{tabular}{l l}
\toprule
\textbf{Kernel} & \textbf{MAPE}\\
\midrule
cudnnConvolutionBackwardFilter & 9.16\% \\
elementwise\_kernel & 24.35\% \\
CatArrayBatchedCopy\_aligned16\_contig & 14.96\% \\
Memset & 34.27\% \\
triton & 4.13\% \\
cudnnConvolutionBackwardData & 7.89\% \\
tensor\_kernel\_scan\_innermost\_dim & 153.91\% \\
MemcpyDtoH & 37.05\% \\
cublasSgemm\_v2 & 37.08\% \\
softmax\_warp\_forward & 229.09\% \\
MemcpyHtoD & 27.71\% \\
cudnnConvolutionForward & 6.31\% \\
multi\_tensor\_apply\_kernel & 1.51\% \\
cublasSgemmStridedBatched & 63.61\% \\
nll\_loss\_backward\_reduce\_cuda\_kernel\_2d & 253.28\% \\
softmax\_warp\_backward & 164.03\% \\
unrolled\_elementwise\_kernel & 10.98\% \\
max\_pool\_backward\_nhwc & 17.19\% \\
cublasLtMatmul & 83.92\% \\
MemcpyDtoD & 65.84\% \\
CatArrayBatchedCopy & 96.63\% \\
vectorized\_elementwise\_kernel & 18.18\% \\
distribution\_elementwise\_grid\_stride\_kernel & 228.62\% \\
nll\_loss\_forward\_reduce\_cuda\_kernel\_2d & 171.61\% \\
\bottomrule
\end{tabular}
\caption{Mean absolute percentage error on a held-out validation set, trained on A40 kernel runtimes. Important kernel types for vision models include \texttt{cudnnConvolution} and \texttt{triton}, where we have <10\% prediction error. Kernels with large percentage-wise errors are extremely short in duration, and thus do not impact end-to-end latency significantly.}
\label{tab:a40_kernel_accuracy}
\end{table}




\section{Performance of alternate search algorithms}

Tune supports several search algorithms out of the box. We investigated the performance of a subset of these algorithms by the progress of the search at distinct phases. Each phase was defined by the number of unique valid configurations sampled by the algorithm up to that point. Every algorithm (with the exception of grid search) was allocated a budget of 2000 samples. \Cref{fig:config-search-alg-comparison} shows the results of this experiment, where the MFU is computed from the iteration times predicted by \sysname. Interestingly, despite the fact that these algorithms are general-purpose and therefore lack domain-specific knowledge of the search space, they appear to converge after having sampled about 200 to 300 unique valid configurations, a 60-75\% improvement over grid search.

\section{Fidelity-preserving Tactics for the Megatron-LM Search Space}
\label{appendix:tactics}

We leverage the performance characteristics of certain Megatron-LM configuration knobs to devise four fidelity-preserving tactics, summarized in \Cref{table:fidelity-preserving-tactics}.
\begin{figure}[h!]
    \centering
    \vspace{2em}
\includegraphics[width=0.44\textwidth,keepaspectratio]{figures/alg_comparison_mini.pdf}

\caption{Comparison of search algorithms exploring GPT3-2.7B (left) and GPT3-18.4B (right). Each algorithm is given a 2000 sample budget. Most algorithms achieve near-optimal MFU after 200-300 valid configurations, providing 60-75\% improvement over grid search. \amey{is the model same in both cases?}}
\label{fig:config-search-alg-comparison}
\end{figure}


\begin{table*}[t]
\centering
\begin{tabular}{c p{0.3\textwidth} p{0.4\textwidth}}
\toprule
\textbf{Knob} & \textbf{Performance characteristics} & \textbf{Tactic} \\
\midrule
Activation recomputation & Reduces memory footprint through smart activation checkpointing 
 & If a prior config OOMed with activation recomputation enabled, then skip the similar config that only disables activation recomputation and mark its result as OOMed \\ 
Sequence parallelism & Reduces memory footprint by reducing activation memory with no added communication cost & If a prior config OOMed with sequence parallelism enabled, then skip the similar config that only disables sequence parallelism and mark its result as OOMed \\
Distributed optimizer & Reduces memory footprint by sharding gradient and optimizer state with added communication cost & If a prior config did not OOM without the distributed optimizer, then skip the similar config that only enables the distributed optimizer and set its runtime to be the same \\
No. of microbatches & In the absence of pipeline parallelism, hardware utilization is inversely proportional to the number of microbatches \cite{megatron}. & If a prior config did not OOM with number of microbatches $n$ and no pipeline parallelism, then skip the similar config that only increases the number of microbatches and set its runtime to be the same \\
\bottomrule
\end{tabular}
\caption{Summary of fidelity-preserving tactics used in Megatron-LM configuration search experiment}
\label{table:fidelity-preserving-tactics}
\end{table*}

